% ============================================================================
% Chapter 12 --- Observability and audit
% ----------------------------------------------------------------------------
% Purpose : TAC-correlated event schema, ingestion pipeline, retention,
%           SIEM integration. Belongs to the Control and Governance plane.
%           One additional section (12.2 event schema) compared to the
%           standard layer template (D6.5 extended).
% Page target: ~5 pages.
% Dependencies: Ch. 5 (OpenTelemetry, CloudEvents, syslog RFC 5424),
%               Ch. 7 (D7.6 makes this chapter the canonical schema owner),
%               Ch. 6 / 8 / 9 / 10 / 11 (all layers feed this layer),
%               Ch. 25 (Wave 7 audit/IR builds on this layer).
% Mirrors layer template, plus dedicated event-schema section.
% ----------------------------------------------------------------------------
% Decisions registered in _coherence-EN.md:
%   D12.1 minimum compliance set for the observability layer
%   D12.2 canonical TAC event schema (frozen, referenced everywhere)
%   D12.3 reference instantiation = Vector + Kafka + OpenSearch + Grafana Loki
%   D12.4 alternatives = Splunk + OTel collectors; Sentinel + Log Analytics
%   D12.5 observability sovereignty scores
%   D12.6 retention is policy-driven, not product-driven
% ============================================================================

\chapter{Observability and audit}
\label{ch:observability}

\begin{caveatbox}[Dependencies]
\noindent This chapter applies the layer template of
\trchap{ch:identity}{} to the observability layer, with one
additional section (\S\ref{sec:ob-schema}) that defines the
canonical TAC event schema referenced from \trchap{ch:policy}{}
(D7.6) and from every other layer chapter that produces audit
events. The chapter's outputs are consumed by
\trchap{ch:wave7-audit-ir}{Wave~7 --- TAC-aware audit and IR}. The
layer belongs to the \emph{Control and governance plane}.
\end{caveatbox}

The observability layer is the place where the activity of every
other layer becomes visible to the institution. It is also the
place where two structural risks accumulate quietly: a per-volume
ingestion cost that grows faster than the institution's monitoring
capacity, and a schema drift between sources that makes correlation
unreliable precisely when it is most needed --- during an incident.
The architectural commitment of this layer is consequently to fix
the schema at the institutional level and to keep the ingestion
pipeline portable.

%-----------------------------------------------------------------------------
\section{Functional role and interface contract}
\label{sec:ob-contract}

The observability layer is responsible for: collecting audit-grade
events from every other layer in a uniform schema; transporting
them through a portable ingestion pipeline; retaining them
according to data-classification-driven policies; making them
searchable for routine operations and forensic investigation; and
producing alerting on policy-relevant patterns.

\begin{contractbox}[Observability layer]
\noindent A compliant observability layer, in the sense of this
architecture, must:

\begin{itemize}[leftmargin=1.5em,nosep]
  \item collect events from every layer in the canonical TAC schema
    of \S\ref{sec:ob-schema};
  \item transport events through
    OpenTelemetry-compatible~\cite{opentelemetry-spec} collectors
    and an open intermediary (Kafka~\cite{kafka}, NATS or
    equivalent), so that the analytical store can be substituted
    without re-instrumenting the sources;
  \item retain events under a documented retention schedule keyed
    on data classification, with the schedule reviewable and
    contractually enforceable;
  \item provide search and correlation capabilities sufficient for
    routine operations and for the incident-response runbooks of
    \trchap{ch:wave7-audit-ir}{};
  \item emit alerts on policy-relevant patterns in a form
    consumable by the institution's chosen incident-response
    workflow.
\end{itemize}
\end{contractbox}

\paragraph{Retention is policy-driven, not product-driven.}
Retention durations are determined by the regulatory regime, the
classification of the events, and the institution's incident
response and audit obligations --- not by the storage tier of the
chosen product. The architectural posture is to express retention
as a documented institutional policy and to instantiate it in the
chosen analytical store, never the inverse.

%-----------------------------------------------------------------------------
\section{Audit-store integrity}
\label{sec:ob-integrity}

The audit trail is load-bearing for accountability and for the lawful
basis of monitoring, so the store that holds it must itself be
tamper-evident. The architectural posture requires three controls,
independent of the chosen product:
\begin{itemize}[nosep]
  \item \textbf{Append-only / WORM.} Audit events are written to an
    append-only or write-once-read-many cold tier; operational
    mutation or deletion of past records is not a supported path.
  \item \textbf{Tamper-evidence.} Per-record or per-batch cryptographic
    signing or hash-chaining lets any later alteration or gap be
    detected.
  \item \textbf{Separation of duties.} The operators who run the
    pipeline, the analysts who read the logs, and the authority to
    expire records under the retention schedule are distinct roles;
    retention expiry follows a four-eyes rule.
\end{itemize}
This discharges \gls{gdpr} Art.~32 (integrity of processing) and
Art.~5(1)(f), \gls{nis2} Art.~21(2)(b)/(f) (incident handling;
effectiveness assessment), ISO/IEC~27001
controls A.8.15/A.8.16, and NIST SP~800-53 AU-9 (protection of audit
information). The break-glass procedure does not bypass it: emergency
access is itself an audited, signed event.

%-----------------------------------------------------------------------------
\section{Canonical TAC event schema}
\label{sec:ob-schema}

The canonical TAC event schema is the institutional contract that
every event-producing layer satisfies. Its purpose is to make
correlation across layers possible and to make the audit trail
reconstructable. Table~\ref{tab:ob-schema} lists the required and
recommended fields. The schema is presented as an abstract structure;
its concrete representation in an institutional deployment may use
OpenTelemetry, CloudEvents, or a JSON schema, provided that the
abstract structure is preserved.

\begin{table}[H]
\centering\small
\caption{Canonical TAC event schema. \textbf{R} required for every
event; \textbf{c} required when the field is meaningful for the
producing layer; \textbf{o} optional but recommended.}
\label{tab:ob-schema}
\begin{tabularx}{\textwidth}{%
  >{\hsize=0.8\hsize\linewidth=\hsize\bfseries\raggedright\arraybackslash}X%
  >{\hsize=0.2\hsize\linewidth=\hsize\centering\arraybackslash}X%
  >{\hsize=2.0\hsize\linewidth=\hsize\raggedright\arraybackslash}X%
}
\toprule
Field & Req. & Description \\
\midrule
\texttt{timestamp}              & R & Event time, UTC, ISO~8601, microsecond precision when available. \\
\texttt{correlation\_id}        & R & TAC session identifier or activity correlation identifier; ties events across layers. \\
\texttt{event\_id}              & R & Unique identifier of this event. \\
\texttt{layer}                  & R & Producing layer (identity, policy, network, compute, storage, endpoint, crypto, orchestration). \\
\texttt{event\_type}            & R & Layer-specific event type (e.g.\ \texttt{auth.success}, \texttt{policy.decision}, \texttt{storage.read}). \\
\texttt{subject\_id}            & R & Institutional identifier of the subject involved, where applicable; otherwise \texttt{system}. \\
\texttt{subject\_attributes}    & c & Role activated, group memberships, posture summary at the moment of the event. \\
\texttt{resource\_id}           & c & Identifier of the resource accessed (workload, dataset, endpoint, configuration object). \\
\texttt{resource\_classification} & c & Data classification label of the resource, where applicable. \\
\texttt{environment}            & c & Zone, time-of-day, location signals, posture flags. \\
\texttt{decision\_verdict}      & c & For policy-relevant events: \texttt{allow}, \texttt{deny}, \texttt{indeterminate}. \\
\texttt{decision\_obligations}  & c & Obligations applied as part of the decision (timeout, additional factor, audit class). \\
\texttt{policy\_version}        & c & Version of the policy bundle that produced the decision. \\
\texttt{duration\_ms}           & c & Evaluation or processing duration, milliseconds. \\
\texttt{source}                 & R & Producing component, instance, version. \\
\texttt{severity}               & R & \texttt{info}, \texttt{notice}, \texttt{warn}, \texttt{error}. \\
\texttt{message}                & o & Free-text description for human readers. \\
\bottomrule
\end{tabularx}
\end{table}

\paragraph{Correlation identifier.} The \texttt{correlation\_id} is
the field that turns a stream of events into a navigable activity
trace. It is set when a TAC is activated and propagated to every
event emitted under that activation, across layers, until the TAC
is decommissioned. The mechanism for propagating the identifier is
implementation-specific (HTTP header, message attribute, log field);
the architectural requirement is that propagation occur and that
the identifier be preserved in the event store.

\paragraph{Schema evolution.} The schema as published is the
institutional baseline. Extensions are admissible provided they do
not displace the baseline fields and provided they are documented
in the institution's observability governance. Schema breakage ---
removing a baseline field, redefining its semantics --- is a
governance event, not a product-driven one.

%-----------------------------------------------------------------------------
\section{Reference instantiation: Vector + Kafka + OpenSearch + Grafana Loki + Grafana}
\label{sec:ob-reference}

\begin{instantiationbox}[Observability layer]
\noindent\textbf{Topology.} Vector~\cite{vector-project} collectors
are deployed at every event-producing source and convert local
formats to the canonical TAC schema; Kafka~\cite{kafka} serves as
the durable ingestion buffer; events are partitioned by
classification and routed to two stores:
OpenSearch~\cite{opensearch} for indexed search and structured
analytics, and Grafana Loki~\cite{grafana-loki} for high-volume log
lines whose retention is shorter. Grafana~\cite{grafana} dashboards
present operational and audit views;
alerting flows to the institution's chosen incident-response
workflow. Cold-tier retention beyond the indexed window is held in
S3-compatible object storage (consistent with
\S\ref{sec:st-reference}).
\end{instantiationbox}

\paragraph{Operational considerations.} The reference instantiation
is operationally substantial: the durable ingest buffer, the
indexed-search cluster, and the cold-tier export pipeline each carry
their own operational rhythm. An indicative envelope, for a
research-intensive institution operating the full pipeline, would be
the institution's observability staffing profile sized for the full
pipeline and shared with the security operations
function. The labour market is broad for the individual components
and narrower for the integrated pipeline.

\paragraph{Volume management.} The principal recurring effort is
volume management: not every event is equally valuable, and a
naive ingestion strategy floods the analytical store with noise
that obscures the signal. The architectural posture is to define
event-class retention and indexing tiers as part of the schema
governance, and to enforce them at the collector or buffer rather
than at the analytical store. This reverses the per-GB economy of
the commercial alternatives.

\begin{realworldbox}[Audit pipeline precedents in academia]
\noindent Public documentation of academic SIEM and observability
pipelines is thinner than for identity or compute. Among what is
available in the research-and-education networking community and its
CSIRT fora, several research-intensive
institutions are reported to operate observability pipelines that
combine open-source collectors with commercial analytical stores;
the specific architectures should be confirmed against their own
published material when available.
\end{realworldbox}

%-----------------------------------------------------------------------------
\section{Alternative~1 --- Splunk Enterprise with OpenTelemetry collectors upstream}
\label{sec:ob-alt1}

\begin{alternativebox}[Splunk + OTel]
\noindent The institution operates Splunk Enterprise~\cite{splunk} as the
analytical store and SIEM, with OpenTelemetry collectors as the
upstream collection layer that feeds Splunk while also feeding,
optionally, a parallel open-source store for portability. The
canonical schema of \S\ref{sec:ob-schema} is enforced at the
collector layer, not at Splunk's ingestion.
\end{alternativebox}

\paragraph{Where it adds value.} Splunk's analytical depth, alerting
sophistication, and the depth of its skills market are substantial.
Forensic investigation and complex correlation are well supported.
For institutions with mature Splunk operations, the productivity
gain over a self-built pipeline can be significant.

\paragraph{Where it constrains the institution.} Per-GB ingestion
pricing is the principal financial concern, and one that has been
the subject of considerable institutional revisiting in recent
years. The architectural mitigation is to keep the collectors
portable, so that data flow can be adjusted (sampled, tiered,
re-routed) without re-instrumenting the sources, and so that exit
from Splunk would be a re-routing of the collector layer rather
than a re-foundation of the audit pipeline.

\paragraph{When it could be the right choice.} An institution with
mature Splunk operations, a financial framework that absorbs the
per-GB exposure, and a sovereignty posture that accepts the data
and institutional dimensions trade-offs in exchange for the
analytical depth.

%-----------------------------------------------------------------------------
\section{Alternative~2 --- Microsoft Sentinel with Log Analytics}
\label{sec:ob-alt2}

\begin{alternativebox}[Microsoft Sentinel]
\noindent The institution operates Microsoft Sentinel~\cite{microsoft-sentinel} on Log
Analytics in a designated regional tenancy. Integration with
Entra~ID-based identity and Microsoft 365 estate is direct;
integration with non-Microsoft layers (Linux compute, OpenStack,
Ceph) flows through OpenTelemetry collectors that translate to
the Log Analytics ingest format.
\end{alternativebox}

\paragraph{Where it adds value.} Tight integration with the
Microsoft estate, where one is already in place, and the
cloud-native operating model that reduces
the institutional infrastructure burden.

\paragraph{Where it constrains the institution.} The data
dimension considerations of \trchap{ch:sovereignty-grid}{} apply
here as in \S\ref{sec:nw-sase}: Sentinel processes the institution's
audit events in the supplier's cloud, under conditions that the
institution should evaluate explicitly against its disclosure
regime. The technical and institutional dimensions are bounded by
Sentinel's product roadmap.

\paragraph{When it could be the right choice.} An institution
already operating its identity and endpoint layers on Microsoft
products, where Sentinel's integration depth produces a meaningful
operational simplification, and where the data-dimension trade-off
has been evaluated explicitly.

%-----------------------------------------------------------------------------
\section{Sovereignty grid applied}
\label{sec:ob-sovereignty}

\begin{table}[H]
\centering\small
\caption{Per-dimension sovereignty profile of the observability
instantiations.}
\label{tab:ob-sovereignty}
\begin{tabularx}{\textwidth}{%
  >{\hsize=0.5\hsize\linewidth=\hsize\bfseries\raggedright\arraybackslash}X%
  >{\hsize=1.2\hsize\linewidth=\hsize\centering\arraybackslash}X%
  >{\hsize=1.1\hsize\linewidth=\hsize\centering\arraybackslash}X%
  >{\hsize=1.2\hsize\linewidth=\hsize\centering\arraybackslash}X%
}
\toprule
Dimension & Reference (Vector+Kafka+OpenSearch+Loki) & Splunk + OTel & Sentinel + Log Analytics \\
\midrule
Data          & 3 Robust      & 2 Established & 1 Partial      \\
Technical     & 3 Robust      & 1--2 Partial  & 1 Partial      \\
Financial     & 2 Established & 1 Partial     & 1--2 Partial   \\
Operational   & 2 Established & 3 Robust      & 3 Robust       \\
Institutional & 3 Robust      & 1 Partial     & 1 Partial      \\
\bottomrule
\end{tabularx}
\end{table}

\paragraph{Driving indicators by dimension.} Data: storage location,
disclosure regime, retention policy. Technical: source visibility,
schema preservation across the pipeline. Financial: cost
predictability under volume growth, exit cost. Operational: skills
market depth and analytical sophistication. Institutional: roadmap
control over the schema and over the alerting model.

\noindent The observability layer surfaces the financial dimension
of the sovereignty grid more sharply than most others, because the
per-volume pricing of the commercial alternatives interacts directly
with the institution's instrumentation discipline: the more visible
the institution makes itself to itself, the more it pays the
supplier. The architectural mitigation is to decouple instrumentation
from analytical-store choice, so that the institution can manage
volume tactically without losing the activity it has chosen to
make visible.

%-----------------------------------------------------------------------------
\section{Regulatory anchoring}
\label{sec:ob-regulatory}

\noindent The observability layer produces the records on which
several regulatory obligations rest. Under the applicable data
protection regime it
contributes to its records-of-processing obligation (derivable from
the canonical TAC schema), its breach-notification obligations
(through detection and timeline reconstruction) and its
security-of-processing obligation (through monitoring effectiveness).
Under the applicable cybersecurity baseline regime it
operationalises the incident-handling requirement
and the regime's incident-reporting timelines (early warning,
full notification, and final report), with the audit
pipeline as the primary timeline-reconstruction substrate. The
layer is the principal anchor of 27001~A.8.15 (logging) and
A.8.16 (monitoring activities) and contributes to NIST CSF
\textsf{DETECT} (DE.CM, DE.AE) and \textsf{RESPOND} (RS.AN ---
analysis). Detailed mapping in \trchap{ch:regulatory-mapping}{}
\S\ref{sec:rm-gdpr}, \S\ref{sec:rm-nis2}, \S\ref{sec:rm-iso},
\S\ref{sec:rm-nist-csf}.

%-----------------------------------------------------------------------------
\section{Common pitfalls}
\label{sec:ob-pitfalls}

\begin{pitfallbox}[Per-GB cost trap]
\noindent A commercial SIEM whose pricing scales with ingestion
volume creates an incentive to under-instrument: events are
sampled, dropped, or routed elsewhere to keep the bill manageable,
and the visibility of the institution to itself becomes a financial
variable rather than a security one. The mitigation is the
collector-layer decoupling described in
\S\ref{sec:ob-reference}: ingestion volume becomes a tactical
question that does not require re-instrumenting the sources.
\end{pitfallbox}

\begin{pitfallbox}[Schema drift between sources]
\noindent Each producing layer expresses events in its own format;
without a canonical schema enforced at the collector, the
analytical store accumulates events whose fields cannot be
reliably correlated. Incident response that depends on
\texttt{correlation\_id} discovers, mid-incident, that the field
is present only on some sources. The mitigation is to make the
canonical schema of \S\ref{sec:ob-schema} a contractual obligation
of the collector layer, and to verify the obligation in the
conformance suite of \trchap{ch:conformance}{}.
\end{pitfallbox}

\begin{pitfallbox}[Undocumented retention]
\noindent The retention duration of the analytical store is the
result of the chosen storage tier, the cleanup job's schedule, and
the supplier's defaults --- which together produce an actual
retention different from the institution's intended policy. The
mitigation is to publish the retention schedule per event class
and to verify it through deliberate retrieval at the documented
horizon.
\end{pitfallbox}

\begin{pitfallbox}[Sampling that breaks correlation]
\noindent Volume management often takes the form of sampling: keep
one event in $N$ for a given class. Sampling that drops events from
the middle of a TAC activation breaks correlation: the surviving
events of the activation lose their context. The mitigation is to
sample by TAC session, not by event: either an entire session is
retained, or none is. Per-event sampling should be reserved for
event classes that do not participate in correlation.
\end{pitfallbox}
